\subsection{Product Perspective}

\subsubsection{Class Diagrams}

\subsubsection{State Diagrams}

\subsubsection{Scenarios}

\textbf{User Creates an Account}\\
J. Doe is a frequent cyclist, and desires an application that gives them the opportunity to log their cycling 
activity, and find information about the best bike routes in his city. They find that BBP perfectly 
suits their needs, and proceed to install the application. The application at first only shows the 
path viewing functionality: J. Doe wants to also record their trips, so they proceed to the ‘sign up’ 
section. The system then asks them to set a username, email and password, which J. Doe inserts, along 
with some optional personal information that they choose not to include. Inserting the username ‘JD’, 
the system responds that this username has already been taken, and prompts to insert a new username. 
Upon insertion of a different username, which was not taken by any other user, the system notifies 
that an email was sent to the specified address with a code, which is required to complete the 
registration. Upon inserting the code, the system creates the new account, and J. Doe can begin 
using the complete application. \\
\newline
\textbf{Registered User Automatically Records a Trip}\\
J. Doe is about to go biking in their city. Before going out to get their bike, they decide that they 
want to record this trip using BBP, so they open the app, where they are logged in to their account, 
and initiate an automated recording of a new trip. The system starts acquiring GPS and sensor data 
from the smartphone, and waits until J. Doe starts biking: at that point, the speed detected by the 
system is high enough to start tracking the device in order to reconstruct the path. While biking, 
J. Doe encounters a pothole: their bike lurches to one side, and the system perceives from the mobile 
device’s sensors that an abnormal movement has occurred, therefore it notes that an obstacle was 
encountered in that location. At the end of their trip, J. Doe gets off their bike and stops the 
recording, at which BBP notifies that the trip was successfully saved to their trip record. \\
\newline
\textbf{Registered User Manually Records a Trip}\\
J.Doe has finished cycling along a new route they found in their city, and they want to save the path 
in BBP while it is still fresh in their memory, since they didn’t activate the automated recording 
before starting to bike. They go to the ‘manual trip insertion’ section of the application, which 
requests the insertion of a starting and ending location and time, and the names and status of all 
roads on the route, in the order in which they were traversed. Additionally the system gives the 
option to mark any obstacles encountered along the route: J. Doe remembers seeing some potholes on 
some of the roads, so they mark them. Once all the information has been inserted, they press ‘save trip’, 
at which point the system reconstructs the path from the provided data and saves the trip to J. Doe’s 
record. \\
\newline
\textbf{User Looks at Possible Bike Paths}\\
J. Doe wants to bike to a new attraction that has recently opened in their city. However, he has never 
been to that neighbourhood, so they want to see if they can find any paths that can quickly and easily 
take them there. They open the BBP and go to the ‘path viewing’ section: after inserting their current 
location as an origin, and the attraction’s location as a destination, the system searches for any published 
paths which are suitable for this route. It finds a couple of paths, and calculates a score for each 
one before displaying them on a map: the one with the highest score is highlighted. J. Doe selects the 
highlighted path and checks what additional information is provided: BBP shows that the path isn’t the 
most direct way to reach the destination, but it is relatively flat and devoid of obstacles. J. Doe then 
selects the next highest scoring path: this path is more direct, but has more obstacles and roads in 
worse conditions, thus scoring lower. J. Doe considers that they want to reach the destination as quickly 
as possible, and thinks that their mountain bike is sturdy enough to not care too much about the worse 
road conditions, therefore they press the ‘export path’ button and choose to export the path to Google 
Maps, so they may begin navigation.

\subsection{Product Functions}

\textbf{Sign up \& Login} \\
Users can access the app in two ways: 
\begin{itemize}
    \item Unregistered access: accessing the app without registering an account gives a user access to the 
    path visualization interface. By inserting a start and end location, the application will show on a map 
    all the paths between the two points, sorted based on their score (as described in detail in the 'Path 
    Visualization' product function).
    \item Registered access: the app also gives users the option to register a personal account. It will ask 
    the user to insert a username, e-mail and a password, along with some optional personal data (name, date 
    of birth, profile picture). If the username or email is already taken by another user, the system will 
    notify that a different username must be chosen. A verification email will be sent to the provided address 
    with an OTP, which will be requested by the system before the registration can be submitted. Once the 
    registration is complete, the user can login to their account, which gives them the ability to record, and 
    eventually publish, their bike trips, in addition to the features of the unregistered access. The user can 
    log out from their account at any time. \\
\end{itemize}
\textbf{Recording of Trips} \\
The main functionality of BBP is to allow registered users to record and save their bike trips.
There are two ways to record a trip in BBP:
\begin{itemize}
    \item Manual recording: if the user wants to record their trip manually, they must insert all relevant 
    information into the BBP application: the start and end locations of their trip, the names of the streets 
    present in the path (in the order in which they were traversed), their status (selected from predefined 
    status values: optimal, medium, sufficient, requires maintenance), and the start and end times of the trip. 
    Users can also manually mark the locations of obstacles encountered along the path. The system will 
    reconstruct the path based on the streets provided and their order.
    \item Automatic recording: if the user wants to record their trip automatically, they must begin the 
    automated recording before the start of their trip, and allow BBP to acquire data from their mobile device. 
    BBP will then start tracking the GPS position of the user's mobile device to determine what path they are 
    taking; the tracking will begin once the speed of the mobile device is estimated to be elevated enough to 
    assume that they have started biking (e.g. 15 km/h), and will not end until the users manually halts the 
    recording. BBP will simultaneously access some of the mobile device's sensors, such as the gyroscope or 
    the accelerometer, so that if a notable movement is registered (due to the presence of some obstacle on 
    the road, such as a pothole) an obstacle marker is added to that point in the path. The system will also 
    determine the status of the road based on the number of obstacles and sensor data patterns. Once the trip 
    has ended the user must stop the recording, at which point BBP will save the trip.
\end{itemize}
The recorded trips are saved in a database accessible only to the user; trips can be published at any time so 
that other users may visualize them, however if the trip data was recorded automatically BBP will first ask 
the user to verify that the collected information is accurate, and if it is not to manually correct it by 
removing incorrect obstacle markers and adjusting the status of the roads, before publishing the trip. \\
\newline
\noindent
\textbf{Trip Statistics Generation} \\
Once a user has recorded a trip, the system will automatically generate some statistics about the trip. The 
statistics include:
\begin{itemize}
    \item total distance covered
    \item average speed
    \item time taken to complete the trip
    \item total elevation difference between start and end points
    \item effective elevation traversed
    \item average slope of the path 
\end{itemize}
These statistics are calculated based on the recorded path data and the inserted start/end times. The system 
will also query an external service to obtain meteorological data relative to the path at that time: if it is 
available, it will be added to the path statistics. The computed statistics are stored with the corresponding 
path in the database; the users have the option to access them when visualizing a path.
Additionally, the system will compute aggregate statistics across all the trips the user has recorded (e.g. 
average time taken to complete trips, average elevation difference of trips, etc.) \\
\newline
\textbf{Trip History Management} \\
Users have the option to view all their past recorded trips. The trips are presented in a list format in 
chronological order, with the date of recording, the name and the start and end locations immediately visible. 
By selecting a specific trip, the user can view the computed statistics for that trip, and a map view of the 
path: there is an option to open the path in a third party navigation app (e.g. Google Maps), so that the user 
can navigate the same route again if they desire. 
The user has the option to edit or delete the trips at any time, and can also choose to publish the trip: if 
the data for the trip was automatically generated and the user has not yet reviewed it,  the system will 
prompt the user to verify and correct the information before publishing. If a trip that has already been 
published is deleted or modified, the system will propagate the action to the published trip (e.g. if a user 
deletes a published trip from their personal trip history, the trip will no longer be available to other users); 
the user also has the option to make a published trip private again at any time, immediately removing it from 
public availability. \\
\newline
\textbf{Path Visualization} \\
Users, both registered and unregistered, can visualize published bike trips on a map. The application will 
prompt the user to insert a start and end location: once inserted, the system will search for all paths that 
connect those two points. The system then computes a score for each of the paths: the score is based on the 
status of the roads in the path, the time it would take for the user to reach the destination by following 
the path (calculated by considering the total length of the path and the elevation difference), the amount of 
obstacles along the path and the difficulty of the path (calculated by considering the average slope of the 
path). The system then shows the user all the identified paths on a map: the path with the highest score will 
be highlighted, while the others will be shown with lower opacity, and the scores will all be visible. Only 
the 5 paths with the highest score will be immediately visible, to avoid cluttering; the user will have the 
option to make all possible paths visible. The user can select a path by clicking on it: the application will 
then give the user the option to view the statistics of the path, select the path as the main highlighted path 
(if it isn't already), and to open the path in a third party navigation app (e.g. Google Maps). The user will 
also have the option to view all paths in a list, sorted by score.

\subsection{User Characteristics}
\subsubsection*{Unregistered User}
An unregistered user can use BBP in a reading-only way, visualizing paths on the map and their score to receive suggestions on which bike trip to take. 
Furthermore at every moment he can decide to register to take advantage of all the features of BBP.

\subsubsection*{Registered User}
A registered user is the main character of the app, and is fundamental for its correct functioning, because the database of routes of BPP is based on insertion of registered users. 
In fact, a registered user can visualize paths on the map like an unregistered user and they can also insert routes, manually or automatically, and make them public to let all users see them.


\subsection{Assumptions, Dependencies and Constraints}
\subsubsection{Domain Assumptions}
\begin{itemize}
    \item D1: Each registered user has a valid email address
    \item D2: Users have a mobile device capable of running the BBP application
    \item D3: Users' mobile devices have GPS capabilities, with typical smartphone precision (5-10 meters)
    \item D4: GPS signals are generally available with sufficient precision while cycling
    \item D5: Registered users who use the automated recording functionality have accelerometer and gyroscope sensors in their mobile devices
    \item D6: Users' have internet access when using online features (e.g. path publication and visualization, weather data retrieval)
    \item D7: Users bike on roads intended for cycling (where bike tracks exist, or cars are rare and speed limits are compatible with biking)
    \item D8: Users who manually record trips bike on named roads which are marked on maps
    \item D9: A sufficient number of registered users are willing to publish their bike routes
    \item D10: Obstacles on bike paths cause a significant enough shift in movement to be detectable through the mobile device's sensors
    \item D11: Users can assess road surface conditions and categorize them in the predefined status values
    \item D12: Bike paths can be reconstructed reasonably accurately based on the names of the streets and the order in which they were traversed
\end{itemize}

\subsubsection{Dependencies}
\subsubsection*{Mapping and routing service}
The BBP system depends on an external mapping and routing service (e.g GoogleMaps, OpenStreetMap, etc.) for tracking and visualizing the bike paths.
\subsubsection*{Meteorological service}
The BBP system depends on an external, third-party meteorological service to acquire data relative to the weather conditions on the bike paths, which are then used to enrich the data about the path.
